% Zone „Kennst du schon" – Daten, Kl. 7
\setcounter{aufgabe}{0}
\zonekopf

\begin{aufgabe}{\textbf{Ich kann eine Strichliste auszählen und angeben, welcher Anteil das ist.} Die Kinder der 7c haben ihr Lieblingsobst genannt.}

\sachtabelle{lc}{Obst & Strichliste}{Apfel & \strichliste{7}\\ Banane & \strichliste{5}\\ Erdbeere & \strichliste{9}\\ Kiwi & \strichliste{3}}

\begin{teile}
\teil Wie viele Kinder mögen Apfel am liebsten? \leerfeld
\teil Welches Obst wurde am häufigsten genannt? \leerfeld
\teil Wie viele Kinder wurden insgesamt befragt? \leerfeld
\teil Welcher Anteil aller Kinder mag Erdbeere am liebsten?\\ Bruch: \leerfeld \quad Prozent: \leerfeld[\%]
\teil Welcher Anteil aller Kinder mag Kiwi am liebsten?\\ Bruch: \leerfeld \quad Prozent: \leerfeld[\%]
\end{teile}
\end{aufgabe}

\begin{aufgabe}{\textbf{Ich kann ausrechnen, wie viel Prozent ein Teil vom Ganzen ist.} Gib den Anteil in Prozent an.}
\begin{teilezwei}
\tz 50 von 100: \leerfeld[\%] & \tz 3 von 10: \leerfeld[\%] \\
\tz 18 von 40: \leerfeld[\%] & \\
\end{teilezwei}
\begin{teile}
\teil Schreibe $5\,\%$ als Dezimalzahl. \leerfeld
\teil $45\,\%$ der Kinder kommen mit dem Rad, $30\,\%$ mit dem Bus, der Rest zu Fuß. Wie viel Prozent kommen zu Fuß?\\ \leerfeld[\%]
\teil In einer Gruppe sind 14 Mädchen und 6 Jungen. Wie viel Prozent der Gruppe sind Mädchen?\\ \leerfeld[\%]
\end{teile}
\end{aufgabe}

\begin{aufgabe}{\textbf{Ich kann Werte aus einem Säulendiagramm ablesen.} Das Diagramm zeigt, wie viele Besucher ein Museum in den ersten fünf Monaten hatte.}

\saeulenab[ymax=60,ystep=10,yfein=5,ylabel=Besucher in Tausend]{Januar/40,Februar/25,März/35,April/50,Mai/45}

Lies ab, wie viele Tausend Besucher kamen.
\begin{teilezwei}
\tz Januar: \leerfeld Tausend & \tz April: \leerfeld Tausend \\
\tz Februar: \leerfeld Tausend & \\
\end{teilezwei}
\begin{teile}
\teil Wie viele Besucher waren es im März? Schreibe die ganze Zahl. \leerfeld
\end{teile}
\end{aufgabe}

\begin{aufgabe}{\textbf{Ich kann Zahlen ordnen und Werte an einer Skala ablesen.} Ordne die Zahlen der Größe nach, mit der kleinsten beginnend.}
\begin{teile}
\teil $7;\ 3;\ 9;\ 5$\\ \leerfeld\,$<$\,\leerfeld\,$<$\,\leerfeld\,$<$\,\leerfeld
\teil $21;\ 12;\ 18;\ 15$\\ \leerfeld\,$<$\,\leerfeld\,$<$\,\leerfeld\,$<$\,\leerfeld
\teil $1{,}7;\ 1{,}65;\ 1{,}8;\ 1{,}72$\\ \leerfeld\,$<$\,\leerfeld\,$<$\,\leerfeld\,$<$\,\leerfeld
\end{teile}
Welche Zahlen gehören zu den Punkten A und B?

\zahlenstrahl[xmin=1.5,xmax=2,xstep=0.05,karo=1.3]{1.65/A, 1.85/B}

\begin{teilezwei}
\tz A: \leerfeld & \tz B: \leerfeld \\
\end{teilezwei}
\end{aufgabe}

\begin{aufgabe}{\textbf{Ich kann Längen in Zentimetern und Millimetern angeben und zeichnen.} Rechne um.}
\begin{teilezwei}
\tz $3\,\text{cm} =$ \leerfeld[mm] & \tz $60\,\text{mm} =$ \leerfeld[cm] \\
\tz $4{,}5\,\text{cm} =$ \leerfeld[mm] & \tz $8\,\text{mm} =$ \leerfeld[cm] \\
\end{teilezwei}
\begin{teile}
\teil Zeichne eine Strecke von $7{,}5\,\text{cm}$ Länge.\par\vspace{1.2cm}
\end{teile}
\end{aufgabe}

\begin{aufgabe}{\textbf{Ich kann Winkel als Teil des Vollkreises sehen und mit dem Geodreieck zeichnen.} Ein Vollkreis hat $360^\circ$. Gib den Teil des Vollkreises an.}
\begin{teile}
\teil $180^\circ$ sind \leerfeld\ des Vollkreises.
\teil $90^\circ$ sind \leerfeld\ des Vollkreises.
\end{teile}
Wie viel Grad sind es?
\begin{teile}
\teil ein Drittel des Vollkreises: \leerfeld[$^\circ$]
\teil ein Sechstel des Vollkreises: \leerfeld[$^\circ$]
\end{teile}
\end{aufgabe}

\begin{aufgabe}{\textbf{Ich kann Winkel als Teil des Vollkreises sehen und mit dem Geodreieck zeichnen – weiter.} Trage den Winkel am Scheitel S an.}
\noindent\begin{minipage}[t]{0.32\linewidth}
\begin{teile}
\teil $40^\circ$
\end{teile}
\winkelstrahl[4.6]
\end{minipage}\hfill
\begin{minipage}[t]{0.32\linewidth}
\begin{teile}
\teil $75^\circ$
\end{teile}
\winkelstrahl[4.6]
\end{minipage}\hfill
\begin{minipage}[t]{0.32\linewidth}
\begin{teile}
\teil $130^\circ$
\end{teile}
\winkelstrahl[4.6]
\end{minipage}
\end{aufgabe}

\begin{aufgabe}{\textbf{Ich kann mit Dezimalzahlen rechnen und sinnvoll runden.} Rechne im Kopf.}
\begin{teilezwei}
\tz $2{,}5 + 1{,}5 =$ \leerfeld & \tz $8{,}4 : 2 =$ \leerfeld \\
\end{teilezwei}
Rechne mit dem Taschenrechner.
\begin{teilezwei}
\tz $0{,}4 \cdot 250 =$ \leerfeld & \\
\end{teilezwei}
Runde auf eine Nachkommastelle.
\begin{teilezwei}
\tz $8{,}46 \approx$ \leerfeld & \tz $17{,}96 \approx$ \leerfeld \\
\tz $23 : 7 \approx$ \leerfeld & \\
\end{teilezwei}
\end{aufgabe}

\begin{aufgabe}{\textbf{Ich kann den kleinsten und den größten Wert, die Spannweite, den Mittelwert und den Median einer Liste bestimmen.} Fünf Kinder haben beim Dartwerfen diese Punkte erzielt: $14;\ 9;\ 17;\ 11;\ 20$.}
\begin{teilezwei}
\tz kleinster Wert: \leerfeld & \tz größter Wert: \leerfeld \\
\tz Spannweite: \leerfeld & \tz Mittelwert: \leerfeld \\
\tz Median: \leerfeld & \\
\end{teilezwei}
\end{aufgabe}

\begin{aufgabe}{\textbf{Ich kann ausrechnen, was „die Hälfte“, „ein Drittel mehr“ oder „doppelt so viel“ bedeutet.} Rechne.}
\begin{teilezwei}
\tz die Hälfte von 80: \leerfeld & \tz doppelt so viel wie 35: \leerfeld \\
\tz ein Drittel von 240: \leerfeld & \tz ein Drittel mehr als 150: \leerfeld \\
\tz ein Viertel weniger als 120: \leerfeld & \\
\end{teilezwei}
\end{aufgabe}

\begin{aufgabe}{\textbf{Ich finde den Fehler, wenn ein Anteil vom falschen Ganzen berechnet wurde.} In einem Chor singen 15 Mädchen und 10 Jungen. Ben rechnet so:}

\rechnung{\text{Anteil der Jungen} &= 10 : 15 \\ &\approx 0{,}67 = 67\,\%}

\begin{teile}
\teil Finde den Fehler, benenne ihn und korrigiere die Rechnung.
\teil In einer Kiste liegen 21 rote und 9 grüne Äpfel. Wie viel Prozent der Äpfel sind grün?\\ \leerfeld[\%]
\end{teile}
\end{aufgabe}
